% !TeX spellcheck = en_US
\documentclass[a4paper,12pt]{article}
\usepackage{listings}
%\usepackage[T1]{fontenc}
\usepackage{fontspec}
\usepackage[greek, english]{babel}
\usepackage[super]{nth}
\usepackage{fancyhdr}
\usepackage[left=1.50cm, right=1.50cm, top=2cm, bottom=1.cm, includeheadfoot]{geometry}
\usepackage{amsmath}
\usepackage{graphicx}
\usepackage{subcaption}
\usepackage{float}
%\usepackage[colorinlistoftodos]{todonotes}
\usepackage{xcolor}
\usepackage{ulem}
\usepackage{amsmath}
\usepackage{caption}
\usepackage{subcaption}
\usepackage{tabularray}
\usepackage[framed]{matlab-prettifier}
\usepackage{svg}
\usepackage{nth}
\usepackage{wrapfig}
\usepackage{siunitx}
\sisetup{output-exponent-marker=\ensuremath{\mathrm{e}}}

\usepackage[pdfauthor={Alexandra Gianni, Nikos Stylianou},
pdftitle={Second Problem Set},
pdfcreator={TeX},
pdfsubject={ECE447 - Neuro-Fuzzy Computing}]{hyperref}


\hypersetup{
	colorlinks,
	citecolor=blue,
	filecolor=black,
	linkcolor=black,
	urlcolor=blue
}

\lstset{
	style              = Matlab-editor,
	basicstyle         = \footnotesize,
	escapechar         = ",
	mlshowsectionrules = true,
}

\setmainfont{Times New Roman}

\pagestyle{fancy}
\fancyhf{}
\fancyhead[l]{\footnotesize \nth{2} Problem Set}
\fancyhead[r]{\footnotesize Neuro-Fuzzy Computing}
%\fancyfoot[r]{\footnotesize \thepage}
\renewcommand{\footrulewidth}{0.4pt} % Line at the footer visible

\fancyfoot[c]{\thepage}

\fancypagestyle{first}{
	\fancyhf{}
	\renewcommand{\headrulewidth}{0pt}
	\fancyfoot[c]{{\large \today}}
	\renewcommand{\footrulewidth}{0.0pt} % Line at the footer visible
} 

\newcommand{\MathSpace}{\ }
\setlength{\parindent}{0pt} % Paragraph identation is set to 0

\newcommand\ddfrac[2]{\frac{\displaystyle #1}{\displaystyle #2}} % Better fractions for inline

\renewcommand{\thesubsection}{(\alph{subsection})} % Change subsection style
\renewcommand{\thesection}{\hspace{-18pt}} % Make section style look nicer
\renewcommand\thesubsubsection{\thesection\hspace{18pt}(\alph{subsection}). \roman{subsubsection}}

\definecolor{veraman}{HTML}{55928A}

\begin{document}
	
	\begin{titlepage}
		\thispagestyle{first}
		
		\newcommand{\HRule}{\rule{\linewidth}{0.5mm}} % Defines a new command for the horizontal lines, change thickness here
		
		\center % Center everything on the page
		
		%----------------------------------------------------------------------------------------
		%	HEADING SECTIONS
		%----------------------------------------------------------------------------------------
		
		\textsc{\LARGE University of Thessaly}\\[1.6cm] % Name of your university/college
		\includegraphics[scale=.5]{Images/uth-logo.png}\\[1cm] % Include a department/university logo - this will require the graphicx package
		\textsc{\Large Neuro-Fuzzy Computing}\\[0.6cm] % Major heading such as course name
		\textsc{\large ECE447}\\[0.5cm] % Minor heading such as course title
		
		%----------------------------------------------------------------------------------------
		%	TITLE SECTION
		%----------------------------------------------------------------------------------------
		
		\HRule \\[0.5cm]
		{ \huge \nth{2} Problem Set}\\[0.4cm] % Title of your document
		\HRule \\[1.8cm]
		
		%----------------------------------------------------------------------------------------
		%	AUTHOR SECTION
		%----------------------------------------------------------------------------------------
		
		
		\vspace*{1cm}
		\begin{minipage}{\textwidth}
			\centering
			\begin{tblr}{cc}
				 \emph{{\LARGE Alexandra Gianni}} & \emph{{\LARGE Nikos Stylianou}} \\ [3mm]
				 \emph{{\LARGE ID: 3382}} & \emph{{\LARGE ID: 2917}} \\
			\end{tblr}
		\end{minipage}\\[2.5cm]
		
		% If you don't want a supervisor, uncomment the two lines below and remove the section above
		%\Large \emph{Author:}\\
		%John \textsc{Smith}\\[3cm] % Your name
		
		%----------------------------------------------------------------------------------------
		%	DATE SECTION
		%----------------------------------------------------------------------------------------
		
		
		%\vfill % Fill the rest of the page with whitespace
		
	\end{titlepage}
	
	% !TeX spellcheck = en_US
\section{Problem 1}
In this exercise we need to find the minimum of the given 2-dimensional function:\\
\begin{equation}	
		F(\mathbf{w})=w_{1}^{2}+w_{2}^{2}+(0.5w_{1}+w_{2})^{2}+(0.5w_{1}+w_{2})^{4}
	\label{eq:function1}	
\end{equation}
with the Conjugate Gradient (Fletcher-Reeves) method and the Gradient Descent.\\

Initially, we can conclude that the function $F(w)$ is not in quadratic form because of the term $(0.5w_{1}+w_{2})^{4}$.
A function is said to be in quadratic form if it can be expressed as a second-degree polynomial where all the terms are either squared terms or cross-products of the variables. The presence of the fourth-degree term $(0.5w_{1}+w_{2})^{4}$.
makes this function a higher-degree polynomial, specifically a quartic function with respect to $(0.5w_{1}+w_{2})$, which means it cannot be classified as quadratic.\\
Also, the independent values in this function are $w_{1},w_{2}$, because only with them we can manipulate the $F(w)$.\\ 

As an initial guess we have $w\left(0\right) = \left[3, 3\right]^T$.\\
Τhe steps we have to use are specific for each iteration

\begin{center}
	\underline{\textit{FIRST ITERATION k = 0}}
\end{center}

\underline{Step1: Calculate the Gradient at  $w\left(k\right)$ }\\
\(\nabla f(w_1,w_2) = \left(\begin{array}{c}
	\dfrac{\partial f}{\partial w_1} \\[4mm]
	\dfrac{\partial f}{\partial w_2}
\end{array}\right)\) = $\left(\begin{array}{c}
	2w_1 + (0.5w_1+w_2) + 2(0.5w_1+w_2)^3\\[1mm]
	2w_2 + 2(0.5w_1+w_2) + 4(0.5w_1+w_2)^3
\end{array}\right) = \left(\begin{array}{c}
	2.5w_1 + w_2 + 2(0.5w_1+w_2)^3\\[1mm]
	w_1 + 4w_2 + 4(0.5w_1+w_2)^3
\end{array}\right)$ \\[3mm]

where at the point $w\left(0\right) = \left[3, 3\right]^T$ we have $\nabla f(x) = \left(\begin{array}{c}
	-53 \\
	-19
\end{array}\right)$
\\[4mm]

\begin{center}
	\underline{\textbf{GRADIENT DESCENT}}
\end{center}
Gradient descent is one of the most favored optimization technique, because of its simplicity and its generality.\\
In order to find the minimum of the Function ~\ref{eq:function1} we must first initialize our initial point $w_{0}$, value of tolerance $tol$ and Step size\\
Given the same initial guess $w\left(0\right) = \left[3, 3\right]^T$. \\
In order to find the minimum of the function~\ref{eq:function1} we need to follow the Gradient Descent's steps:\\

\underline{Step 1: Initialize initial point $w_{0}$, value of tolerance $tol$ and Step size}\\
As we have mentioned, the initial values are:
\begin{itemize}
	\item $w\left(0\right) = \left[3, 3\right]^T$
	\item $tol = 10^{-6}$
	\item $\alpha = 0.01$\\
	The exercise says that we apply the GD method with unit step movement. However, if we set $\alpha = 1$, which is considered a large value, the algorithm may not converge to the optimal point and it will even diverge completely. That's why we set $\alpha = 0.01$. 
\end{itemize} 
\vspace{2mm}

Taking these initializations into account, we follow the Gradient Descent's steps:\\

\underline{Step 1: Compute gradient $\nabla f$ at $w_{k}$}
\vspace{4mm}

\underline{Step 2: Make a scaled step in the opposite direction to the gradient}\\
\begin{equation}
	\begin{gathered}
		\text{step} = \alpha \cdot \nabla f_{k}\\
	\end{gathered}
\end{equation}
\vspace{2mm}

\underline{Step 3: Check convergence}\\
If the step size is smaller than the convergence, then we stop. Else, we repeat Step 1 and Step 2 while we update $x$.
\vspace{4mm}

\underline{Step 4: Update $x$}
\begin{equation}
	x_{k+1} = x_{k} - \alpha \cdot \nabla f_{k}
\end{equation}
\vspace{4mm}

Taking these steps into consideration we will attempt to minimize our function.\\

\begin{center}
	\underline{\textit{ITERATION k = 0}}
\end{center}

\underline{Step 1: Compute Gradient}\\
\(\nabla f(w_1,w_2) = \left(\begin{array}{c}
	\dfrac{\partial f}{\partial w_1} \\[4mm]
	\dfrac{\partial f}{\partial w_2}
\end{array}\right)\) = $\left(\begin{array}{c}
	2w_1 + (0.5w_1+w_2) + 2(0.5w_1+w_2)^3\\[1mm]
	2w_2 + 2(0.5w_1+w_2) + 4(0.5w_1+w_2)^3
\end{array}\right) = \left(\begin{array}{c}
	2.5w_1 + w_2 + 2(0.5w_1+w_2)^3\\[1mm]
	w_1 + 4w_2 + 4(0.5w_1+w_2)^3
\end{array}\right)$ \\[3mm]

where at the point $w\left(0\right) = \left[3, 3\right]^T$ we have $\nabla f_{0} = \left(\begin{array}{c}
	192.75 \\
	379.5
\end{array}\right)$
\\[4mm]

\underline{Step 2: Step size}
\[
step = a \cdot \nabla f_{0} = 0.01 \cdot \left(\begin{array}{c}
	192.75 \\
	379.5
\end{array}\right) \rightarrow step =\left(\begin{array}{c}
1.92750 \\
3.7950
\end{array}\right)
\]
\\[4mm]

\underline{Step 3: Check Convergence}\\
The norm of the step size is $\| step \| = 4.256440 > tol$, so we update $w$
\\[4mm]

\underline{Step 4: Update}
\[ 
w_1 = w_0 - a\cdot \nabla f_0 =  \left[\begin{array}{c}
	3 \\
	3
\end{array}\right] - \left(\begin{array}{c}
1.92750 \\
3.7950
\end{array}\right) \rightarrow
w_1 = \left[\begin{array}{c}
	 1.073\\
	 -0.795
\end{array}\right]
\]
\\[4mm]

\begin{center}
	\underline{\textit{ITERATION k = 1}}
\end{center}

\underline{Step 1: Compute Gradient}\\
\(\nabla f(w_1,w_2) = \left(\begin{array}{c}
	\dfrac{\partial f}{\partial w_1} \\[4mm]
	\dfrac{\partial f}{\partial w_2}
\end{array}\right)\) $= \left(\begin{array}{c}
	2.5w_1 + w_2 + 2(0.5w_1+w_2)^3\\[1mm]
	w_1 + 4w_2 + 4(0.5w_1+w_2)^3
\end{array}\right)$ \\[3mm]

where at the point $w\left(1\right) = \left[1.073, -0.795\right]^T$ we have $\nabla f_{1} = \left(\begin{array}{c}
	1.851603 \\
	-2.176795
\end{array}\right)$
\\[4mm]

\underline{Step 2: Step size}
\[
step = a \cdot \nabla f_{1} = 0.01 \cdot \left(\begin{array}{c}
1.851603 \\
-2.176795
\end{array}\right) \rightarrow step =\left(\begin{array}{c}
	0.018516 \\
	-0.021768
\end{array}\right)
\]
\\[4mm]

\underline{Step 3: Check Convergence}\\
The norm of the step size is $\| step \| = \num{2.857773e-02 }> tol$, so we update $w$
\\[4mm]

\underline{Step 4: Update}
\[ 
w_2 = w_1 - a\cdot \nabla f_1 =  \left[\begin{array}{c}
 1.073\\
-0.795
\end{array}\right] - \left(\begin{array}{c}
	0.018516 \\
	-0.021768
\end{array}\right) \rightarrow
w_2 = \left[\begin{array}{c}
	1.054\\
	-0.773
\end{array}\right]
\]
\\[4mm]

\begin{center}
	\underline{\textit{ITERATION k = 2}}
\end{center}

\underline{Step 1: Compute Gradient}\\
\(\nabla f(w_1,w_2) = \left(\begin{array}{c}
	\dfrac{\partial f}{\partial w_1} \\[4mm]
	\dfrac{\partial f}{\partial w_2}
\end{array}\right)\) $= \left(\begin{array}{c}
	2.5w_1 + w_2 + 2(0.5w_1+w_2)^3\\[1mm]
	w_1 + 4w_2 + 4(0.5w_1+w_2)^3
\end{array}\right)$ \\[3mm]

where at the point $w\left(2\right) = \left[1.054, -0.773\right]^T$ we have $\nabla f_{2} = \left(\begin{array}{c}
	1.831867 \\
	-2.098666
\end{array}\right)$
\\[4mm]

\underline{Step 2: Step size}
\[
step = a \cdot \nabla f_{2} = 0.01 \cdot \left(\begin{array}{c}
	1.831867 \\
   -2.098666
\end{array}\right) \rightarrow step =\left(\begin{array}{c}
	0.018319 \\
	-0.020987
\end{array}\right)
\]
\\[4mm]

\underline{Step 3: Check Convergence}\\
The norm of the step size is $\| step \| = \num{2.785702e-02} > tol$, so we update $w$
\\[4mm]

\underline{Step 4: Update}
\[ 
w_3 = w_2 - a\cdot \nabla f_2 =  \left[\begin{array}{c}
	1.054 \\
	-0.773
\end{array}\right] - \left(\begin{array}{c}
	0.018319 \\
-0.020987
\end{array}\right) \rightarrow
w_3 = \left[\begin{array}{c}
	1.036\\
	-0.752
\end{array}\right]
\]
\\[4mm]

\begin{center}
	\underline{\textit{ITERATION k = 3}}
\end{center}

\underline{Step 1: Compute Gradient}\\
\(\nabla f(w_1,w_2) = \left(\begin{array}{c}
	\dfrac{\partial f}{\partial w_1} \\[4mm]
	\dfrac{\partial f}{\partial w_2}
\end{array}\right)\) $= \left(\begin{array}{c}
	2.5w_1 + w_2 + 2(0.5w_1+w_2)^3\\[1mm]
	w_1 + 4w_2 + 4(0.5w_1+w_2)^3
\end{array}\right)$ \\[3mm]

where at the point $w\left(3\right) = \left[1.036, -0.752\right]^T$ we have $\nabla f_{3} = \left(\begin{array}{c}
	1.811156 \\
	-2.024840
\end{array}\right)$
\\[4mm]

\underline{Step 2: Step size}
\[
step = a \cdot \nabla f_{3} = 0.01 \cdot \left(\begin{array}{c}
	1.811156 \\
-2.024840
\end{array}\right) \rightarrow step =\left(\begin{array}{c}
	0.018112 \\
	-0.020248
\end{array}\right)
\]
\\[4mm]

\underline{Step 3: Check Convergence}\\
The norm of the step size is $\| step \| = \num{2.716664e-02} > tol$, so we update $w$
\\[4mm]

\underline{Step 4: Update}
\[ 
w_4 = w_3 - a\cdot \nabla f_3 =  \left[\begin{array}{c}
	1.036\\
	-0.752
\end{array}\right] - \left(\begin{array}{c}
	0.018319 \\
	-0.020987
\end{array}\right) \rightarrow
w_4 = \left[\begin{array}{c}
	1.018\\
	-0.732
\end{array}\right]
\]
\\[4mm]

\begin{center}
	\underline{\textit{ITERATION k = 4}}
\end{center}

\underline{Step 1: Compute Gradient}\\
\(\nabla f(w_1,w_2) = \left(\begin{array}{c}
	\dfrac{\partial f}{\partial w_1} \\[4mm]
	\dfrac{\partial f}{\partial w_2}
\end{array}\right)\) $= \left(\begin{array}{c}
	2.5w_1 + w_2 + 2(0.5w_1+w_2)^3\\[1mm]
	w_1 + 4w_2 + 4(0.5w_1+w_2)^3
\end{array}\right)$ \\[3mm]

where at the point $w\left(4\right) = \left[1.018, -0.732\right]^T$ we have $\nabla f_{4} = \left(\begin{array}{c}
	1.789642 \\
	-1.954924
\end{array}\right)$
\\[4mm]

\underline{Step 2: Step size}
\[
step = a \cdot \nabla f_{4} = 0.01 \cdot \left(\begin{array}{c}
	1.789642 \\
-1.954924
\end{array}\right) \rightarrow step =\left(\begin{array}{c}
	0.017896 \\
	-0.019549
\end{array}\right)
\]
\\[4mm]

\underline{Step 3: Check Convergence}\\
The norm of the step size is $\| step \| = \num{2.650386e-02} > tol$, so we update $w$
\\[4mm]

\underline{Step 4: Update}
\[ 
w_5 = w_4 - a\cdot \nabla f_4 =  \left[\begin{array}{c}
	1.018\\
	-0.732
\end{array}\right] - \left(\begin{array}{c}
	0.017896 \\
-0.019549
\end{array}\right) \rightarrow
w_5 = \left[\begin{array}{c}
	1.000\\
	-0.712
\end{array}\right]
\]
\\[4mm]

\begin{center}
	\underline{\textit{ITERATION k = 5}}
\end{center}

\underline{Step 1: Compute Gradient}\\
\(\nabla f(w_1,w_2) = \left(\begin{array}{c}
	\dfrac{\partial f}{\partial w_1} \\[4mm]
	\dfrac{\partial f}{\partial w_2}
\end{array}\right)\) $= \left(\begin{array}{c}
	2.5w_1 + w_2 + 2(0.5w_1+w_2)^3\\[1mm]
	w_1 + 4w_2 + 4(0.5w_1+w_2)^3
\end{array}\right)$ \\[3mm]

where at the point $w\left(5\right) = \left[1.000, -0.712\right]^T$ we have $\nabla f_{5} = \left(\begin{array}{c}
	1.767472 \\
	-1.888581
\end{array}\right)$
\\[4mm]

\underline{Step 2: Step size}
\[
step = a \cdot \nabla f_{5} = 0.01 \cdot \left(\begin{array}{c}
	1.767472 \\
-1.888581
\end{array}\right) \rightarrow step =\left(\begin{array}{c}
0.017675 \\
-0.018886
\end{array}\right)
\]
\\[4mm]

\underline{Step 3: Check Convergence}\\
The norm of the step size is $\| step \| = \num{2.586638e-02} > tol$, so we update $w$
\\[4mm]

\underline{Step 4: Update}
\[ 
w_6 = w_5 - a\cdot \nabla f_5 =  \left[\begin{array}{c}
	1.000\\
	-0.712
\end{array}\right] - \left(\begin{array}{c}
0.017675 \\
-0.018886
\end{array}\right) \rightarrow
w_6 = \left[\begin{array}{c}
	0.982\\
	-0.694
\end{array}\right]
\]
\\[4mm]

\begin{center}
	\underline{\textit{ITERATION k = 6}}
\end{center}

\underline{Step 1: Compute Gradient}\\
\(\nabla f(w_1,w_2) = \left(\begin{array}{c}
	\dfrac{\partial f}{\partial w_1} \\[4mm]
	\dfrac{\partial f}{\partial w_2}
\end{array}\right)\) $= \left(\begin{array}{c}
	2.5w_1 + w_2 + 2(0.5w_1+w_2)^3\\[1mm]
	w_1 + 4w_2 + 4(0.5w_1+w_2)^3
\end{array}\right)$ \\[3mm]

where at the point $w\left(5\right) = \left[1.000, -0.712\right]^T$ we have $\nabla f_{5} = \left(\begin{array}{c}
	1.767472 \\
	-1.888581
\end{array}\right)$
\\[4mm]

\underline{Step 2: Step size}
\[
step = a \cdot \nabla f_{5} = 0.01 \cdot \left(\begin{array}{c}
	1.767472 \\
	-1.888581
\end{array}\right) \rightarrow step =\left(\begin{array}{c}
	0.017675 \\
	-0.018886
\end{array}\right)
\]
\\[4mm]

\underline{Step 3: Check Convergence}\\
The norm of the step size is $\| step \| = \num{2.586638e-02} > tol$, so we update $w$
\\[4mm]

\underline{Step 4: Update}
\[ 
w_6 = w_5 - a\cdot \nabla f_5 =  \left[\begin{array}{c}
	1.000\\
	-0.712
\end{array}\right] - \left(\begin{array}{c}
	0.017675 \\
	-0.018886
\end{array}\right) \rightarrow
w_6 = \left[\begin{array}{c}
	0.982\\
	-0.694
\end{array}\right]
\]
\\[4mm]

\begin{center}
	\underline{\textit{ITERATION k = 7}}
\end{center}

\underline{Step 1: Compute Gradient}\\
\(\nabla f(w_1,w_2) = \left(\begin{array}{c}
	\dfrac{\partial f}{\partial w_1} \\[4mm]
	\dfrac{\partial f}{\partial w_2}
\end{array}\right)\) $= \left(\begin{array}{c}
	2.5w_1 + w_2 + 2(0.5w_1+w_2)^3\\[1mm]
	w_1 + 4w_2 + 4(0.5w_1+w_2)^3
\end{array}\right)$ \\[3mm]

where at the point $w\left(5\right) = \left[1.000, -0.712\right]^T$ we have $\nabla f_{5} = \left(\begin{array}{c}
	1.767472 \\
	-1.888581
\end{array}\right)$
\\[4mm]

\underline{Step 2: Step size}
\[
step = a \cdot \nabla f_{5} = 0.01 \cdot \left(\begin{array}{c}
	1.767472 \\
	-1.888581
\end{array}\right) \rightarrow step =\left(\begin{array}{c}
	0.017675 \\
	-0.018886
\end{array}\right)
\]
\\[4mm]

\underline{Step 3: Check Convergence}\\
The norm of the step size is $\| step \| = \num{2.586638e-02} > tol$, so we update $w$
\\[4mm]

\underline{Step 4: Update}
\[ 
w_6 = w_5 - a\cdot \nabla f_5 =  \left[\begin{array}{c}
	1.000\\
	-0.712
\end{array}\right] - \left(\begin{array}{c}
	0.017675 \\
	-0.018886
\end{array}\right) \rightarrow
w_6 = \left[\begin{array}{c}
	0.982\\
	-0.694
\end{array}\right]
\]
\\[4mm]








	% !TeX spellcheck = en_US
\section{Problem 3}

For the given neural network, we have:
\begin{itemize}
	\item learning rate $LR = 1$,
	\item $w^1\left(0\right) = -3,\ w^2\left(0\right) = -1$,
	\item $b^1\left(0\right) = 2,\ b^2\left(0\right) = -1$ and
	\item input/target pair $\left\{p=1,\ t=0\right\}$
\end{itemize}
\vspace*{1mm}

\begin{center}
	\underline{\textit{FIRST ITERATION}}
\end{center}

\underline{Step 1: Calculate first layer's output}
\[
\begin{gathered}
n^1 = w^1 p + b^1 = (-3)(1) + 2 = -1\\
a^1 = {Swish}\left(n^1\right) = {Swish}\left(-1\right) = \dfrac{n^1}{1+e^{-n^1}} = \dfrac{-1}{1+e} = -0.2689
\end{gathered}
\]

\underline{Step 2: Calculate second layer's output}
\[
\begin{gathered}
n^2 = w^2 a^1 + b^2 = (-1)(-0.2689) + (-1) = -0.7311 \\ 
a^2 = {LReLU}\left(n^2\right) = {LReLU}\left(-0.7311\right) = -0.000731
\end{gathered}
\]

\underline{Step 3: Calculate error}
\[
e = t-a^2 = \left(0-\left(-0.000731\right)\right) = 0.000731
\]

\underline{Step 4: Calculate sensitivity on second layer}
\[
s^2 = -2\ {LReLU}^{'}\left(n^2\right)\left(t-a^2\right) = -2 \left(0.001\right) \left(0.000731\right) = \num{-1.462e-06}
\]
\textit{\small LReLU's derivative is $1$ for $x>0$ and $0.001$ for $x<0$.}\\ 

\underline{Step 5: Calculate sensitivity on first layer using back-propagation}
\[
\begin{gathered}
s^1 = Swish^{'} \left(n^1\right) \left(w^2\right)^T s^2 = Swish^{'} \left(-1\right) \left(-1\right) (\num{-1.462e-06}) = 0.0723 (-1) (\num{-1.462e-06}) \\
s^1 = \num{1.0570e-07}
\end{gathered}
\]

\underline{Step 6: Update wheights and biases}
\[
\begin{gathered}
	w^{2}(1)=w^{2}(0)-LR\ s^{2}(a^{1})^{T} = -1 - 1(\num{-1.462e-06})(-0.2689) \approx -1 \\
	 b^{2}(1)=b^{2}(0)-LR\ s^{2} = -1 - 1(\num{-1.462e-06}) \approx -1 \\
	 w^{1}(1)=w^{1}(0)-LR\ s^{1}(a^{0})^{T} = -3 - 1(\num{1.0570e-07})(-1) \approx -3 \\
	 b^{1}(1)=b^{1}(0)-LR\ s^{1} = 2 - 1(\num{1.0570e-07}) \approx 2
\end{gathered}
\]

Since there were no changes on the biases and weights, the next iteration will not change the parameters of the given neural network, but we will calculate them anyway.
\vspace*{1mm}
\begin{center}
	\underline{\textit{SECOND ITERATION}}
\end{center}

\underline{Step 1:}
\[
\begin{gathered}
	n^1 = w^1 p + b^1 = (-3)(1) + 2 = -1\\
	a^1 = {Swish}\left(n^1\right) = {Swish}\left(-1\right) = \dfrac{n^1}{1+e^{-n^1}} = \dfrac{-1}{1+e} = -0.2689
\end{gathered}
\]

\underline{Step 2:}
\[
\begin{gathered}
	n^2 = w^2 a^1 + b^2 = (-1)(-0.2689) + (-1) = -0.7311 \\ 
	a^2 = {LReLU}\left(n^2\right) = {LReLU}\left(-0.7311\right) = -0.000731
\end{gathered}
\]

\underline{Step 3:}
\[
e = t-a^2 = \left(0-\left(-0.000731\right)\right) = 0.000731
\]

\underline{Step 4:}
\[
s^2 = -2\ {LReLU}^{'}\left(n^2\right)\left(t-a^2\right) = -2 \left(0.001\right) \left(0.000731\right) = \num{-1.462e-06}
\]

\underline{Step 5:}
\[
\begin{gathered}
	s^1 = Swish^{'} \left(n^1\right) \left(w^2\right)^T s^2 = Swish^{'} \left(-1\right) \left(-1\right) (\num{-1.462e-06}) = 0.0723 (-1) (\num{-1.462e-06}) \\
	s^1 = \num{1.0570e-07}
\end{gathered}
\]

\underline{Step 6:}
\[
\begin{gathered}
	w^{2}(1)=w^{2}(0)-LR\ s^{2}(a^{1})^{T} = -1 - 1(\num{-1.462e-06})(-0.2689) \approx -1 \\
	b^{2}(1)=b^{2}(0)-LR\ s^{2} = -1 - 1(\num{-1.462e-06}) \approx -1 \\
	w^{1}(1)=w^{1}(0)-LR\ s^{1}(a^{0})^{T} = -3 - 1(\num{1.0570e-07})(-1) \approx -3 \\
	b^{1}(1)=b^{1}(0)-LR\ s^{1} = 2 - 1(\num{1.0570e-07}) \approx 2
\end{gathered}
\]
	
	
	
\end{document}